% Flavor catalog per-process template.
% process_id: CR011
% family: collider_rs
% owner_agent_id: PKA-CR011
% writer_agent_id: null
% checker_agent_id: null
% opus_signoff_id: null
% Canonical metadata sidecar: CR011.yaml
% Status is controlled by the YAML sidecar status_history, not this comment.

\section{\texorpdfstring{Longitudinal Vector-Boson Scattering}{Longitudinal Vector-Boson Scattering}}

\subsection*{Process}
CR011 covers the strong-electroweak-symmetry-breaking collider tail probed by
longitudinal vector-boson scattering.  The assignment shorthand is
\(pp\to W_L W_L W_L W_L\); experimentally, the cleanest current promoted
observable is same-sign \(W_L^\pm W_L^\pm\) scattering in association with two
forward jets, \(pp\to jj\,W_L^\pm W_L^\pm\), with leptonic \(W\) decays used
to select the final state.  The process is an EW-precision-tail entry for the
custodial/composite electroweak sector, not a quark-flavor amplitude.  In the
collider-RS catalog it is the longitudinal-boson scattering counterpart to
explicit electroweak KK-resonance searches.

\subsection*{Current best limit(s)}
The canonical current limit promoted in the sidecar is the ATLAS Run-2
longitudinal same-sign \(WW\) result in
\texttt{atlas\_2025\_arxiv2503\_11317.txt}.  ATLAS reports the most stringent
constraint to date on fully longitudinal same-sign \(WW\) production:
\[
  \sigma_{\rm fid}\!\left(pp\to jj\,W_L^\pm W_L^\pm\right)
  < 0.45~\mathrm{fb}\quad(95\%~\mathrm{CL}),
\]
with an expected limit of \(0.70~\mathrm{fb}\).  The same analysis reports
evidence for production with at least one longitudinally polarized \(W\) boson
at \(3.3\sigma\), using \(140~\mathrm{fb}^{-1}\) of \(13~\mathrm{TeV}\) data.
For historical comparison, CMS set
\(\sigma(pp\to jj\,W_L^\pm W_L^\pm)<1.17~\mathrm{fb}\) at \(95\%~\mathrm{CL}\)
in \texttt{cms\_2020\_arxiv2009\_09429.txt}.

The PDG QGC review snapshot
\texttt{pdg2025\_wz\_quartic\_couplings.txt} is used as framework context for
anomalous-quartic-gauge-coupling conventions.  It does not provide a single
combined longitudinal-\(WW\) fiducial limit or RS mass exclusion, so no value
from that PDG review is promoted as the CR011 numerical observable.

\subsection*{Relevance to RS}
In custodial RS and composite-Higgs-like completions, longitudinal VBS tests
the sector that unitarizes electroweak symmetry breaking: electroweak KK
vectors, scalar resonances, and the pattern of couplings to \(W_L\) and \(Z_L\).
This is complementary to colored KK-gluon and vectorlike-partner searches.  A
VBS excess or exclusion constrains the electroweak/composite resonance sector
only after specifying the custodial gauge group, the resonance spectrum, the
branching fractions into vector bosons versus tops or composite fermions, and
the production mechanism.

For the anarchic-flavor pipeline in this repository, CR011 is not expected to
be the leading constraint.  The quark-scan methodology note quotes
\(\Mkk^{\min}(p50,\gs=3)=47.26~\mathrm{TeV}\) from the low-energy
\(\Delta F=2\) scan, while direct electroweak VBS reach remains a few-TeV
collider cross-check unless a non-anarchic or protected flavor structure
weakens the flavor bound.  CR011 is therefore useful as a diagnostic of the
strong-EWSB sector and as leverage for non-anarchic RS variants, not as a
replacement for the active low-energy flavor constraints.

\subsection*{Post-2010 developments}
\texttt{atlas\_2017\_arxiv1611\_02428.txt} established the LHC same-sign
\(WWjj\) VBS/aQGC program at \(8~\mathrm{TeV}\), separating electroweak and
strong production regions and introducing VBS-sensitive anomalous-coupling
fiducial selections.  \texttt{atlas\_2020\_arxiv2004\_10612.txt} extended the
VBS program to rare \(ZZjj\) electroweak production, important because neutral
diboson scattering tests a different electroweak channel and different
background systematics.

\texttt{cms\_2020\_arxiv2009\_09429.txt} was the first CMS polarization
measurement in same-sign \(WWjj\), setting the \(1.17~\mathrm{fb}\) longitudinal
\(WW\) upper limit used above as the historical comparator.
\texttt{atlas\_2024\_arxiv2312\_00420.txt} then provided a full Run-2 ATLAS
same-sign \(WWjj\) measurement and aQGC interpretation, including differential
information and a same-sign \(H^{\pm\pm}\) VBF search.  The current
polarization-specific improvement is the ATLAS 2025 result promoted here.

For EFT/aQGC context, \texttt{atlas\_2026\_arxiv2603\_18630.txt} is the current
ATLAS combined interpretation of measurements sensitive to quartic electroweak
gauge couplings.  Its coefficient intervals are not quoted in this draft
because CR011's promoted observable is the direct longitudinal \(WW\) fiducial
cross-section limit; adding coefficient-level values would require a separate
table extraction and a choice of unitarity/positivity convention.

\subsection*{Validity / model dependence}
The ATLAS/CMS limits are fiducial and polarization-template dependent.  They
assume the Standard Model electroweak VBS topology, leptonic \(W\)-decay event
selection, and fitted longitudinal/transverse polarization templates.  They do
not directly exclude a universal RS \(M_{KK}\) value.  A reinterpretation must
simulate the chosen custodial-RS spectrum and account for interference with the
SM VBS amplitude, resonance widths, acceptance, and decay branching fractions.

For aQGC/EFT statements, the interpretation depends on the operator basis,
whether one or several Wilson coefficients are profiled, and whether unitarity
constraints are imposed.  For explicit RS resonances, the assumptions differ
again: a narrow electroweak vector resonance decaying visibly to \(VV\) is much
more constrained by VBS-like searches than a broad resonance with dominant
decays to top partners, heavy fermions, or other composite-sector states.  This
entry should therefore be treated as collider-sector cross-check information
until a dedicated RS likelihood or recast is attached.

\subsection*{Code coverage}
\textbf{NO}.  A targeted grep for VBS, longitudinal-polarization, diboson, and
aQGC terms across \texttt{quarkConstraints/}, \texttt{qcd/},
\texttt{flavorConstraints/}, \texttt{neutrinos/}, \texttt{yukawa/},
\texttt{warpConfig/}, \texttt{solvers/}, \texttt{scanParams/}, and
\texttt{tests/} returned no matches.  Adjacent code evidence shows KK-scale
bookkeeping and low-energy flavor evaluation, not collider direct-search
filtering: \texttt{quarkConstraints/scan.py:359} maps
\(\Lambda_{\rm IR}\) to \(M_{KK}\), \texttt{quarkConstraints/scan.py:377}
evaluates \(\Delta F=2\) constraints, and \texttt{scanParams/scan.py:528}
passes \(M_{KK}\) into the lepton-flavor-violation check.  No ATLAS/CMS,
VBS, \(WWjj\), \(ZZjj\), polarization, or collider-likelihood module is
present in the searched paths.

\subsection*{Implementation difficulty}
\textbf{HIGH}.  A live CR011 constraint needs more than a scalar threshold.  It
requires a collider reinterpretation stack: event generation for the chosen
custodial-RS resonance spectrum, showering and detector/selection acceptance,
polarization-sensitive templates or an external recast tool such as CheckMATE,
MadAnalysis5, or SModelS where applicable, and a likelihood interface for the
ATLAS/CMS fiducial regions.  EFT/aQGC coefficient comparisons are simpler to
record, but mapping them into a concrete RS mass bound is model dependent and
not implemented in this repository.

\subsection*{Key references}
\texttt{atlas\_2025\_arxiv2503\_11317}; \texttt{cms\_2020\_arxiv2009\_09429};
\texttt{atlas\_2024\_arxiv2312\_00420}; \texttt{atlas\_2020\_arxiv2004\_10612};
\texttt{atlas\_2017\_arxiv1611\_02428}; \texttt{atlas\_2026\_arxiv2603\_18630};
\texttt{pdg2025\_wz\_quartic\_couplings};
\texttt{contino\_2011\_arxiv1109\_1570}.
